% Created 2026-05-10 Sun 16:49
% Intended LaTeX compiler: lualatex
\documentclass[11pt]{article}
\usepackage{amsmath}
\usepackage{fontspec}
\usepackage{graphicx}
\usepackage{longtable}
\usepackage{wrapfig}
\usepackage{rotating}
\usepackage[normalem]{ulem}
\usepackage{capt-of}
\usepackage{hyperref}
\usepackage{fontspec}
\setmainfont{UT Sans}[
BoldFont={UT Sans Bold},
UprightFont={UT Sans Regular},
FontFace={m}{n}{UT Sans Medium}
]
% Fallback fake italic using a slant
\usepackage{etoolbox}
\newcommand{\fakeitalic}[1]{{\addfontfeatures{FakeSlant=0.18}#1}}
\AtBeginDocument{\renewcommand{\emph}[1]{\fakeitalic{#1}}}
\usepackage{minted}
\usepackage[a4paper,margin=3cm]{geometry}
\usepackage{lastpage}
\usepackage{fancyhdr}
\usepackage{lastpage}
\pagestyle{fancy}                % set fancy style
\fancyhf{}                        % clear all headers and footers
\cfoot{\thepage/\pageref*{LastPage}}  % center footer
\renewcommand{\headrulewidth}{0pt}    % remove header line
\renewcommand{\footrulewidth}{0pt}    % remove footer line
% also override plain page style for first page
\fancypagestyle{plain}{%
\fancyhf{}%
\cfoot{\thepage/\pageref*{LastPage}}%
\renewcommand{\headrulewidth}{0pt}%
\renewcommand{\footrulewidth}{0pt}%
}
\usepackage{url}
\author{Elias-Robert BAJCSI}
\date{\textit{<2026-05-10 Sun>}}
\title{Laborator 7\\\medskip
\large PLFP}
\hypersetup{
 pdfauthor={Elias-Robert BAJCSI},
 pdftitle={Laborator 7},
 pdfkeywords={},
 pdfsubject={},
 pdfcreator={Emacs 30.2 (Org mode 9.7.11)}, 
 pdflang={English}}
\begin{document}

\maketitle
\section{\texttt{./Cargo.toml}}
\label{sec:org2db891c}
\begin{minted}[breaklines=true,breakanywhere=true,breaksymbol=,fontsize=\small,linenos=false]{toml}
[package]
name = "lab7"
version = "0.1.0"
edition = "2021"

[[bin]]
name = "ex1"
path = "src/bin/ex1.rs"
[[bin]]
name = "ex2"
path = "src/bin/ex2.rs"
[[bin]]
name = "ex3"
path = "src/bin/ex3.rs"
\end{minted}
\section{\texttt{./src/bin/ex1.rs}}
\label{sec:org05637c9}
\begin{minted}[breaklines=true,breakanywhere=true,breaksymbol=,fontsize=\small,linenos=false]{rs}
// regula 1+2+3+4: depinde de stare globala, fara parametri, nedeterminist daca starea se schimba, efect secundar
static TVA: f64 = 0.21;
fn calculeaza_tva_impura() -> f64 { 100.0 * TVA }

struct Appender {
    sufix: String,
}

impl Appender {
    fn append_to_word_impura(&self, baza: &str) -> String {
        format!("{}{}", baza, self.sufix)
    }
}

fn append_to_word_pura(baza: &str, sufix: &str) -> String {
    format!("{}{}", baza, sufix)
}

fn pret_total(pret: f64, tva: f64) -> f64 {
    pret + pret * tva
}

fn adauga(v: &[i32], item: i32) -> Vec<i32> {
    let mut rez = v.to_vec();
    rez.push(item);
    rez
}

fn main() {
    let a = Appender { sufix: "!".into() };
    println!("{}", calculeaza_tva_impura());
    println!("{} {}", a.append_to_word_impura("Salut"), append_to_word_pura("Salut", "!"));
    println!("{} {:?}", pret_total(100.0, 0.21), adauga(&[1, 2, 3], 4));
}

#[cfg(test)]
mod tests {
    use super::*;

    #[test]
    fn test_append_pura() {
        assert_eq!(append_to_word_pura("ana", "_x"), "ana_x");
    }

    #[test]
    fn test_pret_total_pura() {
        assert_eq!(pret_total(10.0, 0.21), 12.1);
    }

    #[test]
    fn test_adauga_pura() {
        let v = vec![1, 2, 3];
        let v2 = adauga(&v, 4);
        assert_eq!(v, vec![1, 2, 3]);
        assert_eq!(v2, vec![1, 2, 3, 4]);
    }
}
\end{minted}
\section{\texttt{./src/bin/ex2.rs}}
\label{sec:orgdaaa16a}
\begin{minted}[breaklines=true,breakanywhere=true,breaksymbol=,fontsize=\small,linenos=false]{rs}
fn f(x: f64, y: f64) -> Option<f64> {
    let num = x + y;
    let den = (x - y).powi(2);
    if den == 0.0 {
        return None;
    }
    let inside = num / den;
    if inside < 0.0 {
        return None;
    }
    let root = inside.sqrt();
    let div = 4.0 * root;
    if div == 0.0 {
        return None;
    }
    Some(x / div)
}

fn aplica_tva_catalog(produse: &[(String, f64)], tva: f64) -> Vec<(String, f64)> {
    produse.iter().map(|(n, p)| (n.clone(), p + p * tva)).collect()
}

fn compune_pure(f: fn(i32) -> i32, g: fn(i32) -> i32) -> impl Fn(i32) -> i32 {
    move |x| f(g(x))
}

fn dublu(x: i32) -> i32 { x * 2 }
fn adauga3(x: i32) -> i32 { x + 3 }

fn main() {
    println!("{:?}", f(10.0, 2.0));
    let p = vec![("Paine".to_string(), 5.5), ("Cafea".to_string(), 20.0)];
    println!("{:?}", aplica_tva_catalog(&p, 0.21));
    println!("{}", compune_pure(dublu, adauga3)(7));
}

#[cfg(test)]
mod tests {
    use super::*;

    #[test]
    fn test_f_ok() {
        assert!(f(10.0, 2.0).is_some());
    }

    #[test]
    fn test_f_div_zero() {
        assert_eq!(f(2.0, 2.0), None);
    }

    #[test]
    fn test_f_negative_root() {
        assert_eq!(f(-10.0, 1.0), None);
    }

    #[test]
    fn test_catalog_immutable() {
        let p = vec![("X".to_string(), 100.0)];
        let out = aplica_tva_catalog(&p, 0.21);
        assert_eq!(p[0].1, 100.0);
        assert_eq!(out[0].1, 121.0);
    }

    #[test]
    fn test_compune_pure() {
        let c = compune_pure(dublu, adauga3);
        assert_eq!(c(5), 16);
        assert_eq!(c(5), 16);
    }
}
\end{minted}
\section{\texttt{./src/bin/ex3.rs}}
\label{sec:org5405a30}
\begin{minted}[breaklines=true,breakanywhere=true,breaksymbol=,fontsize=\small,linenos=false]{rs}
#[derive(Clone, Debug)]
struct Produs {
    nume: String,
    pret: f64,
    stoc: u32,
}

fn filtreaza_disponibile(produse: &[Produs]) -> Vec<Produs> {
    produse.iter().filter(|p| p.stoc > 0).cloned().collect()
}

fn aplica_discount(produse: &[Produs], procent: f64) -> Vec<Produs> {
    produse
        .iter()
        .map(|p| Produs { nume: p.nume.clone(), pret: p.pret * (1.0 - procent / 100.0), stoc: p.stoc })
        .collect()
}

fn total_stoc(produse: &[Produs]) -> f64 {
    produse.iter().map(|p| p.pret * p.stoc as f64).sum()
}

fn main() {
    let produse = vec![
        Produs { nume: "Paine".into(), pret: 5.5, stoc: 3 },
        Produs { nume: "Lapte".into(), pret: 7.0, stoc: 0 },
        Produs { nume: "Cafea".into(), pret: 20.0, stoc: 2 },
    ];

    let disponibile = filtreaza_disponibile(&produse);
    let reduse = aplica_discount(&disponibile, 10.0);
    println!("{:.2}", total_stoc(&reduse));
}

#[cfg(test)]
mod tests {
    use super::*;

    #[test]
    fn test_filtreaza() {
        let p = vec![
            Produs { nume: "A".into(), pret: 10.0, stoc: 1 },
            Produs { nume: "B".into(), pret: 20.0, stoc: 0 },
        ];
        assert_eq!(filtreaza_disponibile(&p).len(), 1);
    }

    #[test]
    fn test_discount() {
        let p = vec![Produs { nume: "Paine".into(), pret: 100.0, stoc: 1 }];
        let r = aplica_discount(&p, 10.0);
        assert_eq!(r[0].pret, 90.0);
    }

    #[test]
    fn test_total() {
        let p = vec![
            Produs { nume: "Paine".into(), pret: 10.0, stoc: 2 },
            Produs { nume: "Lapte".into(), pret: 5.0, stoc: 3 },
        ];
        assert_eq!(total_stoc(&p), 35.0);
    }
}
\end{minted}
\end{document}
